\chapter{Your own guard checks, and prompts in Scheme}
\label{ch:checks}

Two more ways to shape a session from \code{init.scm}: tighten the guard
with checks of your own, and rewrite what \nodwyn{} says to the model.

\section{\texttt{register-check}: tighten, never loosen}

\begin{scm}
(register-check "no-compost-talk" "deny"
  (lambda (prim args)
    (and (equal? prim "fs-write")
         (string-contains? (format "~a" args) "compost"))))
\end{scm}

A check has a name, a fixed verdict and a predicate. The predicate gets the
primitive's name and its argument list for \emph{every} primitive call in
the session. When it returns true, the check answers its verdict; otherwise
it has no opinion.

\begin{txt}
you> :eval (fs-write "log.txt" "add compost")
⚑ guard deny fs-write: fs-write is irreversible; changes log.txt; user checker no-compost-talk
error: fs-write: denied: user checker no-compost-talk
\end{txt}

The rules are what make checks safe to write:

\begin{itemize}
  \item The verdict is \code{"deny"} or \code{"confirm"}, never
    \code{"allow"}. Since the strictest verdict wins, a check can only make
    the guard stricter.
\begin{txt}
you> :eval (register-check "x" "allow" (lambda (p a) #t))
error: register-check: register-check: verdict must be "deny" or "confirm", got "allow"
\end{txt}
  \item A predicate that raises an error \textbf{fails closed}: it counts as
    a denial, with the error in the reason.
\begin{txt}
⚑ guard deny fs-list: checker "broken-check" errored: car: argument 1: expected pair, got () ()
\end{txt}
  \item Only you can register a check: the primitive is hidden from the
    model and refused to model-originated code and to the loop's strategy.
  \item A check you add with \code{:eval} applies at once. One you add or
    change in \code{init.scm} applies from the next session: a
    \code{(reload)} does not re-register checks, and says so.
\end{itemize}

\begin{careful}
A \code{"confirm"} check is an \emph{ordinary} confirm: a standing approval
waives it, the judge may answer it, and \code{-{}-yolo} answers it yes. On a
write, writes in the folder going ahead are such an approval, so a
\code{"confirm"} check on a write asks only with \code{[approval] writes =
"ask"}. If a rule must hold with nobody watching, make it a \code{"deny"}.
\end{careful}

\subsection{Checks worth having}

\begin{scm}
;; Never let anything in this session talk to the package index.
(register-check "no-installs" "deny"
  (lambda (prim args)
    (and (equal? prim "shell-exec")
         (pair? args)
         (member (car args) '("pip" "pip3" "npm" "gem" "cargo"))
         (let ((rest (cadr args)))
           (and (list? rest) (member "install" rest))))))

;; The generated folder is built, never edited.
(register-check "hands-off-generated" "deny"
  (lambda (prim args)
    (and (member prim '("fs-write" "fs-edit"))
         (pair? args)
         (string-prefix? "gen/" (format "~a" (car args))))))
\end{scm}

\begin{tip}
Write predicates defensively: they see every call, including ones whose
arguments are not what you expect. Check \code{(pair? args)} before
\code{(car args)}, and use \code{(format "\textasciitilde a" x)} before
treating an argument as a string. A predicate that errors denies the call,
which is safe but annoying.
\end{tip}

\section{Every prompt is Scheme}

Everything \nodwyn{} says to a model is defined in
\file{internal/loop/prompts.scm}, which loads before your \code{init.scm}.
So you can redefine any of it:

\begin{center}
\begin{tabularx}{\linewidth}{lX}
\toprule
\code{(system-prompt repo)} & the control model's system prompt; by default the interactive or the unattended one \\
\code{working-rules} & how to work: shell command lines, the project's environment, reading, the check loop\ldots \\
\code{task-header}, \code{plan-header} & the headings under which the window shows each request, at the head of its message, and its plan, as a message of its own \\
\code{(plan-system-prompt repo)}, \code{(plan-draft-prompt request)}\ldots & the planner's prompts (Chapter~\ref{ch:plans}) \\
\code{compact-summary-prompt} & the request for a handover note (Chapter~\ref{ch:contexts}) \\
\code{finish-no-change-message}\ldots & what the evidence loop says to a model that stops early \\
\code{(reading-reminder-message n)}\ldots & what the default loop says to a model whose last \code{n} rounds only read \\
\code{(run-turn-prompt task complete-line turn max-turns)} & each turn's message in a supervised run \\
\bottomrule
\end{tabularx}
\end{center}

\begin{txt}
you> :eval task-header
"Current task (the user's latest request):"
you> :eval (string-length (session-system-prompt))
5237
\end{txt}

\subsection{Extending the system prompt}

The idiom is to keep the default under another name and wrap it. A raw
string (Chapter~\ref{ch:strings}) keeps the prompt readable:

\begin{scm}
(define base-system-prompt system-prompt)     ; keep the default

(define (system-prompt repo)
  (string-append
   (base-system-prompt repo) "\n\n"
   """
   About this project:
   - Tests run with `make test`; a single package with `go test ./internal/x`.
   - Never edit files under gen/: they are generated by `make gen`.
   - Prefer table-driven tests, as the existing ones are.
   """))
\end{scm}

\begin{careful}
The system prompt and the two headers are read \emph{once}, after
\code{init.scm} runs, and then fixed for the session like the tool list;
they are recorded with the session so a replay rebuilds requests exactly.
A \code{system-prompt} that is not a procedure returning a non-empty string
stops the session from starting, which is what you want: a broken prompt
should fail loudly, not silently send nothing.
\end{careful}

\subsection{The other prompts change live}

Every prompt other than the system prompt and the headers is ordinary
Scheme, read when it is used. Change one in \code{init.scm}, then ask the
running session to reload:

\begin{txt}
you> :eval (reload)
("register-check: the guard's checks are fixed for this session; changes apply from the next one")
\end{txt}

\code{reload} requires the current project revision to be trusted, then runs \code{init.scm} again and returns notes about anything
that could not change mid-session: new tools, new checks, the system
prompt. Everything else (strategies, thresholds, the other prompts, tool
bodies) applies from the next turn.

\begin{tip}
A tool's \emph{body} can change on reload; its \emph{schema} (name,
parameters, doc) cannot. So you can fix a bug in a tool mid-session, but
adding a parameter waits for the next one.
\end{tip}

\begin{puzzle}[ 14: a friendlier reviewer]
Write a check \code{"tests-before-commit"} that denies \code{shell-exec} of
\code{git} with \code{commit} anywhere in its arguments. Then think: why is
this check weaker than it looks? (Hint: what else can run a command
line?)\ans{14}
\end{puzzle}
